\subsection{\texorpdfstring{\(L^{p}\) 中的有向集与 \(L^{p}\) 中的递增序列}{L to the p 中的有向集与 L to the p 中的递增序列}}

我们已定义了（§2, No.~6）序关系 \(\widetilde{f} \leqslant \widetilde{g}\) 于集 \(\widetilde{\mathcal{F}}\) （由在 \(X\) 中几乎处处有定义且有限的数值函数的等价类组成）中；赋有这个序关系及其向量空间结构的 \(\widetilde{\mathcal{F}}\) 是一个 Riesz 空间。No.~5 命题 12 的推论表明：若 \(\widetilde{f}\) 与 \(\widetilde{g}\) 是子空间 \(L to the p\) （\(\widetilde{\mathcal{F}}\) 的）的两个元素，则上确界 \(\sup (\widetilde{f}, \widetilde{g})\) （\(\widetilde{f}\) 与 \(\widetilde{g}\) 在 \(\widetilde{\mathcal{F}}\) 中的） （它是诸函数 \(\sup (f, g)\) 中每一个的类，其中 \(f \in \widetilde{f}\) 且 \(g \in \widetilde{g}\) ）属于 \(L to the p\) ；这特别证明了 \(L to the p\) （赋有由 \(\widetilde{\mathcal{F}}\) 诱导的序关系）是一个 Riesz 空间。

命题 13. --- 在 Riesz 空间 \(L^{p}\) （赋有由范数 \(\|f\|_{p}\) 定义的拓扑）中，映射 \(\widetilde{f} \mapsto |f|\) 是一致连续的，且元素 \(f \geqslant 0\) 所成之集是闭的。

命题的第一部分由 No.~5 的命题 11 立得；由于 \(\widetilde{f} \geqslant 0\) 所成之集也就是 \(\widetilde{f}\) （使 \(|\widetilde{f}| = \widetilde{f}\) 的）所成之集，故它是闭的，因为 \(\widetilde{f} \mapsto |\widetilde{f}|\) 是连续映射且 \(\mathbf{L}^p\) 是 Hausdorff 的。

我们于是看到，\(L^{p}\) 上由范数 \(\|f\|_{p}\) 定义的拓扑与 \(L^{p}\) 的序向量空间结构相容（TVS, II, §2, No.~7）。

命题 14. --- 设 H 是 Riesz 空间 \(L^{p}\) 的一个子集，由类 \(\geqslant 0\) 组成且按关系 \(\leqslant\) 有向。为使 H 在 \(L^{p}\) 中有上确界，必须且只须

\[
\sup _ {\widetilde {f} \in \mathrm{H}} \| \widetilde {f} \| _ {p} <   + \infty .
\]

此时 H 在 \(L^{p}\) 中的上确界是 H 的截面滤子的极限（在 Banach 空间 \(L^{p}\) 中）。

条件显然是必要的，因为 \(\widetilde{f} \mapsto \| \widetilde{f} \|_p\) 是诸元素 \(\geqslant 0\) （\(L to the p\) 的）所成之集上的一个递增函数。为看出它是充分的，我们先注意它蕴涵：H 在映射 \(\widetilde{f} \mapsto \| \widetilde{f} \|_p\) 下的像依单调极限定理在 \(\mathbf{R}\) 中有一个极限；于是 H 的截面滤子 \(\mathfrak{F}\) 在这个映射下的像是 \(\mathbf{R}\) 上一个 Cauchy 滤子的基。若我们能证明 \(\mathfrak{F}\) 本身是 \(L to the p\) 上一个 Cauchy 滤子的基，证明便告完成；因为此时 \(\mathfrak{F}\) 将在 \(L to the p\) 中收敛，由于 \(L to the p\) 是完备的（No.~4, 定理 2），而命题将由 TVS, II, §2, No.~7, 命题 18 得出。

为看出 \(\mathfrak{F}\) 是一个 Cauchy 滤子的基，我们将用到下述引理：

引理. --- 若 \(f\) 与 \(g\) 是 \(\mathcal{L}^p\) 中的两个函数，使 \(0 \leqslant f \leqslant g\) ，则

\[
\left(\mathrm{N} _ {p} (g - f)\right) ^ {p} \leqslant \left(\mathrm{N} _ {p} (g)\right) ^ {p} - \left(\mathrm{N} _ {p} (f)\right) ^ {p}.\tag{9}
\]

当 \(f\) 与 \(g\) 是具有紧支集的连续函数时，关系式 (9) 可写成

\[
\int (g - f) ^ {p} d | \mu | \leqslant \int g ^ {p} d | \mu | - \int f ^ {p} d | \mu |
\]

而它则是初等不等式 \((g - f)^{p} \leqslant g^{p} - f^{p}\) （No.~1, 公式 (2)）的推论。为由此过渡到一般情形，只需注意 (9) 的两个成员是 \(L^{p} \times L^{p}\) 上的连续函数，且每个函数 \(f \geqslant 0\) （\(L^{p}\) 中的）都是具有紧支集的连续函数 \(\geqslant 0\) 的序列的极限（对于 p 阶平均收敛），这由映射 \(g \mapsto |g|\) 在 \(L^{p}\) 上的连续性（命题 11）得出。

引理既已建立，按假设，对每个 \(\varepsilon > 0\) 存在 \(\widetilde{f} \in \mathbf{H}\) ，使对每个 \(\widetilde{g} \geqslant \widetilde{f}\) （属于 \(\mathbf{H}\) 的） ，有 \((\|\widetilde{g}\|_p)^p - (\|\widetilde{f}\|_p)^p \leqslant \varepsilon\) ；由此得出 \((\|\widetilde{g} - \widetilde{f}\|_p)^p \leqslant \varepsilon\) ；于是，若 \(\widetilde{g}_1\) 与 \(\widetilde{g}_2\) 是 \(\mathbf{H}\) 中 \(\geqslant \widetilde{f}\) 的两个元素，则 \(\|\widetilde{g}_1 - \widetilde{g}_2\|_p \leqslant 2\varepsilon^{1/p}\) ，这就证明了 \(\mathfrak{F}\) 是 \(\mathrm{L}^p\) 上的一个 Cauchy 滤子基，并完成了命题 14 的证明。

推论 1. --- 若 \(\widetilde{g}\) 是 H 在 \(\mathbf{L}^p\) 中的上确界，则

\[
\| \widetilde {g} \| _ {p} = \lim _ {\widetilde {f} \in \mathrm{H}} \| \widetilde {f} \| _ {p} = \sup _ {\widetilde {f} \in \mathrm{H}} \| \widetilde {f} \| _ {p}.\tag{10}
\]

这由范数 \(\|\widetilde{f}\|_{p}\) 在 \(L^{p}\) 中的连续性以及单调极限定理得出。

推论 2. --- Riesz 空间 \(L^{p}\) 是完全格序的。

\(L^{p}\) 中每个有向集 H （按关系 \(\leqslant\) ），由类 \(\geqslant 0\) 组成且在 \(L^{p}\) 中有上界者，都有上确界：因为若 \(\widetilde{h}\) 是 H 在 \(L^{p}\) 中的一个上界，则 \(\|f\|_{p} \leqslant \|h\|_{p}\) 对所有 \(\widetilde{f} \in H\) 成立，于是命题 14 可以应用。这就证明了推论（Ch. II, §1, No.~3, 命题 1）。

当把命题 14 的结论对 \(\mathcal{L}^p\) 中的函数而不是它们的类来陈述时，它们不再成立。确切地说，若 M 是 \(\mathcal{L}^p\) 的一个子集，由函数 \(\geqslant 0\) 组成、按关系 \(\leqslant\) 有向，且使 \(\sup_{f\in\mathbb{M}}\mathrm{N}_{p}(f)<+\infty\) ，则 M 的上包络 \(g\) 的类并不

必然等同于在 \(\mathrm{L}^p\) 中诸函数 \(f \in \mathbf{M}\) 的类的上确界；特别地，\(g\) 不一定是 \(p\) 次幂可积的，而且即使 \(g \in \mathcal{L}^p\) ，\(\mathrm{N}_p(g)\) 也可以不同于 \(\sup_{f \in \mathbf{M}} \mathrm{N}_p(f)\) （参见 §1, No.~3, 定理 3 后的注记 1）。

尽管如此，我们有下述定理：

定理 5. --- 设 \((f_n)\) 是函数 \(\geqslant 0\) （\(\mathcal{L}^p\) 中的）组成的一个递增序列。为使这序列的上包络 \(f\) 是 \(p\) 次幂可积的，必须且只须 \(\sup_{n} N sub p(f_n) < +\infty\) 。此时序列 \((f_n)\) \(p\) 阶平均收敛于 \(f\) ，且

\[
\mathrm{N} _ {p} (f) = \sup _ {n} \mathrm{N} _ {p} (f _ {n}) = \lim _ {n \rightarrow \infty} \mathrm{N} _ {p} (f _ {n}).\tag{11}
\]

条件显然必要，我们只需证明它是充分的。今若条件满足，则命题 14 表明序列 \((\tilde{f}_n)\) 是 \(L to the p\) 中的一个 Cauchy 序列，因而序列 \((f_n)\) 是 \(\mathcal{L}^p\) 中的一个 Cauchy 序列；由于 \(f_n(x)\) 趋于 \(f(x)\) 对所有 \(x \in X\) 成立，故 \(f\) 是 \(p\) 次幂可积的，且是序列 \((f_n)\) 对于 \(p\) 阶平均收敛拓扑的极限（No.~4, 定理 3 的推论 1）。于是 \(\mathrm{N}_p(f_n)\) 趋于 \(\mathrm{N}_p(f)\) ，因为 \(\mathrm{N}_p\) 是 \(\mathcal{L}^p\) 上的连续函数。

推论 1. --- 设 \((f_n)\) 是函数 \(\geqslant 0\) （\(\mathcal{L}^p\) 中的）组成的一个递减序列；则序列的下包络 \(f\) 属于 \(\mathcal{L}^p\) ，序列 \((f_n)\) \(p\) 阶平均收敛于 \(f\) ，且

\[
\mathrm{N} _ {p} (f) = \lim _ {n \rightarrow \infty} \mathrm{N} _ {p} (f _ {n}) = \inf _ {n} \mathrm{N} _ {p} (f _ {n}).
\]

前两个断言由把定理 5 应用于序列 \(g_{n} = f_{1} - f_{n}\) （它是递增且有上界的）得出；其余部分则是明显的。

推论 2. --- 设 \((f_n)\) 是 \(\mathcal{L}^p\) 中函数的一个序列。为使上包络 \(f\) （序列 \((f_n)\) 的）是 \(p\) 次幂可积的，必须且只须存在函数 \(g \geqslant 0\) 使 \(\int^{*} g^p d|\mu| < +\infty\) 且 \(f_n \leqslant g\) 对所有 \(n\) 成立。

条件显然必要，取 \(g = f^{+}\) 即可。反之，设条件成立，并令 \(g_{n} = \sup_{k\leqslant n}f_{k}\) ；序列 \((g_{n})\) 是递增的，由 \(p\) 次幂可积函数组成（No.~5, 命题 12 的推论）。正函数的递增序列 \(h_n = g_n + g_1^-\) 满足定理 5 的条件，因为 \(\mathrm{N}_p(h_n)\leqslant \mathrm{N}_p(g + g_1^-) < +\infty\) ；于是它的上包络 \(\sup_{n}h_{n}\) 是 \(p\) 次幂可积的，对 \(f = \sup_{n}h_{n} - g_{1}^{-}\) 亦然。

推论 3. --- 设 A 是一个可数集，\(\mathfrak{F}\) 是 A 上具有可数基的一个滤子，\((f_{\alpha})_{\alpha \in \mathrm{A}}\) 是函数 \(\geqslant 0\) （\(\mathcal{L}^p\) 中的）的一个族。假定存在函数 \(g \geqslant 0\) 使 \(\mathrm{N}_p(g) < +\infty\) 且 \(f_{\alpha} \leqslant g\) 对所有 \(\alpha \in \mathrm{A}\) 成立；则函数 \(\limsup_{\mathfrak{F}} f_{\alpha}\) 是 \(p\) 次幂可积的，且

\[
\limsup _ {\mathfrak {F}} \mathrm{N} _ {p} (f _ {\alpha}) \leqslant \mathrm{N} _ {p} (\limsup _ {\mathfrak {F}} f _ {\alpha}).\tag{12}
\]

设 \((\mathbf{A}_n)\) 是 \(\mathfrak{F}\) 的一个递减基，并令 \(g_{n} = \sup_{\alpha \in \mathbf{A}_{n}}f_{\alpha}\) ；由于 \(\mathbf{A}_n\) 是可数集，由推论 2 知 \(g_{n}\) 是 \(p\) 次幂可积的；另一方面，\(\mathrm{N}_p(g_n)\geqslant \sup_{\alpha \in \mathbf{A}_n}\mathrm{N}_p(f_\alpha)\) 。这样，\(\limsup_{\mathfrak{F}}f_{\alpha}\) 是递减序列 \((g_n)\) 的下包络；于是 \(\limsup_{\mathfrak{F}}f_{\alpha}\) 依推论 1 是 \(p\) 次幂可积的，且

\[
\begin{array}{l} \mathrm{N} _ {p} (\limsup _ {\mathfrak {F}} f _ {\alpha}) = \mathrm{N} _ {p} \left(\inf _ {n} g _ {n}\right) = \lim _ {n \to \infty} \mathrm{N} _ {p} (g _ {n}) \\ \geqslant \lim _ {n \to \infty} \left(\sup _ {\alpha \in \mathrm{A} _ {n}} \mathrm{N} _ {p} (f _ {\alpha})\right) = \limsup _ {\mathfrak {F}} \mathrm{N} _ {p} (f _ {\alpha}). \end{array}
\]

